\begin{abstract}
Operating systems (OS) have become increasingly complex. A bug in their codebase could disrupt critical services and compromise the security of all users. Several projects employ formal methods to verify the correctness of OS kernels, thereby eliminating bugs and delivering robust guarantees. 
Page tables implement virtual-to-physical address translation and enforce memory protection for all processes.
However, proving the correctness of OS kernels' paging functions that traverse or update page tables is a significant challenge. To reduce proof costs, prior work verified a structural lifting property that establishes the equivalence between paging functions operating on concrete multi-level page tables and those operating on an abstract map. However, these proofs still require substantial manual effort. We present a push-button approach to scaling the verification of paging functions. The key feature of our method is that it encodes all verification tasks as quantifier-free formulas, enabling us to leverage SMT solvers to automate proofs. Based on this, we introduce the Push-button Paging Function Verification Framework (PPV) and use it to verify multiple implementations that operate on multi-level page tables and support essential features: dynamic page table allocation and various mapping granularities, including huge pages. PPV is integrated into the Spoq system verification framework to automate functional correctness proofs for the OS codebase, reducing the substantial manual proof burden associated with verifying paging functions.
%This work validates the feasibility and efficiency of using an automated, push-button approach to verify commodity paging functions.
\end{abstract}